
\section{Access sequences with initial trees} 
\label{sec:forbidden} 

\iffalse
Recall that a sequence $X = (x_1, \dots, x_m)$ \emph{avoids} a permutation $\sigma: [k] \rightarrow [k]$, if there are no indices $1 \leq i_1 < i_2 < \cdots < i_k \leq m$, such that the sequence $x_{i_1}, \dots, x_{i_k}$ is order-isomorphic to $(\sigma(1), \dots, \sigma(k))$. 

In this section we prove that the cost of {\sc Greedy} on an arbitrary access sequence (not necessarily a permutation) that avoids a permutation of size $k$ is quasilinear. The result holds for arbitrary initial tree, therefore, it also holds for the cost of {\sc Greedy} restricted to an arbitrary contiguous subsequence of the input. Since preorder sequences avoid 231, as a special case we obtain a quasilinear bound for the {\sc Greedy} analogue of the traversal conjecture for splay trees. Furthermore, we show that the maximum independent rectangle lower bound (or equivalently the cost of {\sc SignedGreedy}) is linear on any such access sequence. We prove the following theorem.
\fi

\subsection{Initial trees.}
\label{sec:initial-tree}
In the BST model, the \emph{initial tree} is the state of the BST before the first element is accessed. In the tree-view of BST execution, we know the 
structure of the tree exactly in every time step, and hence the structure of
the initial tree as well. This is less obvious, however, in the geometric
view. In fact, according to the reduction of DHIKP, the
initial tree is always a \emph{split tree }\cite{DHIKP09}, hidden from us
in geometric view. We show how to ``force'' an arbitrary initial tree into the geometric view of BST execution.

First, we observe a simple fact. Given a binary search tree $T$, we
denote the search path of an element $a$ in $T$ as ${\sf path}(a,T)$. It is easy to show that two BSTs $T_{1}$ and $T_{2}$,
containing $[n]$, have the same structure iff ${\sf path}(a,T_{1})={\sf path}(a,T_{2})$
as a set, for all $a\in[n]$.


The stair of an element is, in some sense, the geometric view of the
search path. Indeed, ${\sf path}(a,T)=\{b\mid b$ is touched if $a$ is
accessed in $T\}$. Similarly, ${\sf stair}_t(a)=\{b\mid b$
is touched if $a$ is accessed at time $t\}$. Thus, to model an initial
tree $T$ into the geometric view, %by Proposition \ref{same tree=same paths},
it is enough to enforce that ${\sf stair}_1(a)={\sf path}(a,T)$.

We can achieve this in a straightforward way, by putting stacks of points in each column below the first row, corresponding to the structure of the initial tree $T$. More precisely, in column $x$ we put a stack of height $d-d(x)+1$, where $d(x)$ is the depth of element $x$ in $T$, and $d$ is the maximum depth of an element (Figure \ref{fig:geo init tree}). One can easily show that ${\sf stair_1(a)}={\sf path(a,T)}$ for all $a\in[n]$.

We remark that the points representing the initial tree are not counted in the cost.


\iffalse
\paragraph{Traversal conjecture.}
The traversal conjecture for splay trees, asked in 1985 by Sleator and Tarjan~\cite{ST85} states that the cost of running the splay tree algorithm on an input $X \in {\sf PreOrd}_n$ with arbitrary initial tree $T$ is $O(n)$. Only the special cases when $X$ is a monotone sequence~\cite{tarjan_sequential, pettie_DS}, or when $T$ is the unique BST underlying $X$~\cite{chaudhuri} are known. The analogous conjecture for {\sc Greedy} can be stated as follows. 

\begin{conjecture}[Traversal for {\sc Greedy}] 
For any input sequence $X \in {\sf PreOrd}_n$, and any initial tree $T$, we have that {\sc Greedy}$(X) = O(n)$.   
\end{conjecture} 

In \S\,\ref{sec:warmup} we prove the special case of this conjecture without initial tree. The initial tree in our proof will be the split-tree created by the reduction of DHIKP, independent of the input sequence $X$. In \S\,\ref{sec:forbidden} we prove a quasilinear upper bound for {\sc Greedy} that applies to preorder sequences with arbitrary initial tree.

\begin{theorem}
\label{thm:pseudo}
Let $T$ be an initial tree. 
Let $S = (x_1, \dots, x_m) \in [n]^m$ be an arbitrary access sequence that avoids a permutation $\sigma$ of size $k$. The cost of {\sc Greedy}, starting with initial tree $T$, on $X = \pset_S$ is $O\left(m \cdot 2^{({\alpha(m)})^{f(k)}}\right)$, for a fixed function $f$ (to be specified later), and $\alpha$ the inverse Ackermann function. 
\end{theorem}



\subsection{Preliminaries}

In this section and in \S\,\ref{sec:block_linear} we use tools from the theory of forbidden matrices. In order to apply these tools, it will be fruitful to encode both the input and the output of {\sc Greedy} in the form of a matrix. All matrices considered in this paper have entries from $\{0,1\}$.


Let $X$ be a {\sc MinASS} input, corresponding to an access sequence $S = (x_1,\dots, x_m)$, where $x_i \in [n]$, and let $Y$ be the output of {\sc Greedy} on $X$ (note that $Y \supseteq X$). It is easy to see that $Y$ consists of points with integral coordinates from the set $[n] \times [m]$. %Thus, a natural encoding of the {\sc Greedy} execution is a matrix with $m$ rows and $n$ columns. More precisely, 
Let us define the \emph{access matrix} $M^X$ as a matrix with $m$ rows and $n$ columns, such that $M^X(i,j) = 1$, iff $(i,j) \in X$, and 0 otherwise. Similarly, let us define the \emph{touch matrix} $T^X$ as a matrix with $m$ rows and $n$ columns, such that $T^X(i,j) = 1$, iff $(i,j) \in Y$, and 0 otherwise. We call the number of ones in a matrix $M$ the \emph{weight} of $M$, denoted as $w(M)$. Observe that the cost of \greedy on $X$ is $w(T^X)$.

%%Equivalently, 
Given $I \subseteq [n]$ and $J \subseteq [m]$, a \emph{submatrix} $M|_{I,J}$ is obtained by removing columns in $[n] \setminus I$ and rows in $[m] \setminus J$ from $M$.   
%contains the rows and columns indexed by $I$ and $J$ respectively. 
%%When $I$ and $J$ contain consecutive numbers, i.e. $I = \{r,\ldots, s\}$, we say that $M|_{I,J}$ is a {\em consecutive submatrix} of $M$.\\
%
We define the \emph{tensor product} between matrices as follows: $M \otimes G$ is the matrix obtained by replacing each one-entry of $M$ by a copy of $G$, and replacing each zero-entry of $M$ by an all-zero matrix equal in size with $G$ (for permutation matrices, in the notation of \S\,\ref{subsec:decomp} we have: $M \otimes G = M[G,G,\dots]$).

\subsection{Forbidden Submatrices}
%%For a matrix $M$, we say that $M'$ is a submatrix of $M$ if $M'$ can be obtained by removing a collection of rows and columns of $M$. %Let $M$ and $P$ be matrices with entries in $\{0,1\}$. 

We say that a matrix $M$ \emph{contains} a matrix $P$ if there is a submatrix $M'$ of $M$ that has the same size as $P$, and $M'(i,j)=1$ whenever $P(i,j)=1$. In words, a 1-entry in $P$ has to be matched with a 1-entry in $M$, and a 0-entry in $P$ can be matched to anything. If $M$ does not contain $P$, we say that $M$ is $P$-\emph{free}, or that $P$ is a \emph{forbidden submatrix} of $M$.

A permutation matrix $M_\sigma$ of a permutation $\sigma$ is defined as: $M(i,j) = 1$, if $\sigma(j) = i$, and 0 otherwise. We observe the following obvious connection between permutation-avoidance and forbidden submatrices. Let $M^X$ be the access matrix corresponding to an access sequence $X$ that avoids permutation $P$. Then $M^X$ is $M_P$-free. Permutation matrices have a single one in every row and in every column. Relaxing this condition slightly, a matrix that has a single one in every column is called a \emph{light} matrix.
 

The theory of forbidden submatrices is mostly concerned with questions of the following type: given an $m$-by-$n$ matrix $M$ with entries $\{0,1\}$, how large can $w(M)$ be, assuming that $M$ is $P$-free, for a fixed matrix $P$. Of the many results of this flavor found in the literature, we collect in the following lemma three that are relevant in this paper.

\begin{lemma}[Forbidden Submatrix Bounds~\cite{furedi_hajnal, marcus_tardos, jfox, klazar, keszegh}]
\label{lem: forbidden}
Let $M$ be an $m$-by-$n$ matrix with entries $\{0,1\}$, $m \geq n$. 

\begin{enumerate}
\item (F\"{u}redi, Hajnal~\cite[Cor.\ 2.7]{furedi_hajnal})~~ if $M$ is $\left( \begin{matrix}
  1 & 0 & 0 \\
  0 & 0 & 1 \\
  0 & 1 & 0
 \end{matrix}\right)$-free, then $w(M) \leq 2(m+n)$.
\item (Marcus, Tardos~\cite{marcus_tardos})~~ if $M$ is $P$-free, for a $k$-by-$k$ permutation matrix $P$, then $w(M) \leq m \cdot f(k)$, for a fixed function $f$.
\item (Klazar~\cite{klazar})~~ if $M$ is $P$-free, for a $k$-by-$k$ light matrix $P$, then $w(M) \leq m \cdot 2^{O(({\alpha(m)})^{g(k)})}$, for a fixed function $g$. 
\end{enumerate} 
\end{lemma}

Statement (ii) of Lemma~\ref{lem: forbidden} was proven by Marcus and Tardos~\cite{marcus_tardos}, answering a question by F\"{u}redi and Hajnal~\cite{furedi_hajnal}. They obtain the function $f(k) = 2k^4\binom{\,k^2}{k} = 2^{O(k \log k)}$. Recently Fox~\cite{jfox} obtained the improved bound $f(k) = 2^{O(k)}$.

Statement (iii) of Lemma~\ref{lem: forbidden} follows from results of Klazar~\cite{klazar}. The proof can also be found in the work of Keszegh~\cite{keszegh}, where the function $g(k) = 2k^2 - 4$ is obtained.

\fi
\iffalse

\begin{proof} 
Let us call the three access points $(a,t_a), (b,t_b), (c,t_c)$ where $t_a < t_b < t_c$. 
There are two cases to consider: First if $b \in \tset_a$, then it is impossible that $c \in \tset_a$ because all points on the left are accessed before the right. This means that there is $d \not \in \set{a,b,c}$ where $d = lca(a,c)$, so in particular, $d$ is accessed before both $a$ and $c$. 
Now we reach a contradiction because the element $c$ in the left-subtree of $d$ is accessed after the element $a$ on the right. 

\lk{cursor}

Otherwise, if $b \not \in \tset_a$, there is an element $d = lca(a,b)$, where $a$ and $b$ are in the left- and right-subtrees of $d$ respectively. 
Element $c$ cannot be in $\tset_d$ for otherwise, the accesses would have alternated from $a$ to $b$ and then back to $c$, so this means that $lca(c,d)$ exists, again contradiction because elements on the right subtrees are accessed before those on the left.     
\end{proof} 
\subsection{Complexity of Access Sequence} 
 
\begin{theorem}
Let $X$ be an input sequence of size $n$ that does not contain a permutation pattern $\rho$ of size $k$. Then $Greedy(X) \leq O_k(n)$ where the $O_k$ notation hides $k$ as constant.  
%If $X$ is $P$-free, then $|{\sc Greedy}(X)| < f(k)(m+n)$, in particular, for any constant $k$, we have $|{\sc Greedy}(X)| = \Theta(m+n)$.
\end{theorem}

\lk{Maybe we could give the exact expressions, since it is interesting in the whole $k=o(loglogn)$ range, and discuss in one sentence as special case the $k$ constant case.}

We remark that our technique allows us to prove this theorem even without the permutation assumption of $X$, i.e. if $X$ is an access sequence of length $m$ on $[n]$, then the cost of {\sc Greedy} on $X$ is $O(m+n)$. 

\fi


\iffalse
\subsection{Proof of Theorem~\ref{thm:pseudo}}
\label{pseudoproof}
Let $M^X$ and $T^X$ be the access, respectively touch matrices corresponding to the execution of \greedy on an input sequence $X$ that avoids a permutation $P$.

We mention that in the case when $P = 231$, a relatively simple argument shows that $T^X$ is $M_P$-free. In words, if the input is $231$-free, then the output of \greedy is also $231$-free. The proof of this claim is simple and we leave it as an exercise for the interested reader. Observe that together with Lemma~\ref{lem: forbidden}(i), this yields an alternative proof of Theorem~\ref{thm:greedy-traversal}, with a matching constant.

For a general avoided permutation $P$, we have to work harder. We define the matrix $G = \left( \begin{matrix}
  0 & 1 & 0 \\
  1 & 0 & 1 
 \end{matrix}\right)$, and call it the \emph{forcing gadget}. We make the following claim.
 \begin{lemma}
 \label{lem:forcing1}
 If $X$ avoids $P$, then $T^X$ is $(M_P \otimes G)$-free.
 \end{lemma}

Observe that Lemma~\ref{lem:forcing1} immediately implies Theorem~\ref{thm:pseudo}(i) via Lemma~\ref{lem: forbidden}(iii), and by the observation that if $P$ is a permutation of size $k$, then $M_P \otimes G$ is a $3k$-by-$2k$ light matrix. (To apply Lemma~\ref{lem: forbidden} we can pad $M_P \otimes G$ with zeros to make it a square matrix). Now we prove Lemma~\ref{lem:forcing1}.

The key idea of our proof is the observation that for every occurrence of the forcing gadget $G$ in the \greedy execution, the bounding box around the gadget must have an access point inside. More precisely:

\begin{lemma}
\label{forcing2}
Let $M' = T^{X}|_{I,J}$ be a submatrix of the touch matrix $T^X$, such that $M'$ contains $G$. Let $I', J'$ denote the coordinates of the smallest bounding box around $M'$, i.e.\ $I' = [\min(I), \max(I)]$, and $J' = [\min(J), \max(J)]$ Then $w(M^{X}|_{I',J'}) \geq 1$.
\end{lemma}

\begin{proof}
Let $(x,t)$ and $(y,t)$ be the two entries in $T^X$ matched with the lower left, respectively lower right 1-entry of $G$. Let $t'$ be the coordinate of the topmost row of $M'$, i.e.\ $t' = \max{(J)}$. The top 1-entry of $G$ needs to match an element in the interval $[x+1, y-1]$ touched in the time interval $[t+1,t']$. Such a touched point, however, cannot exist due to Lemma~\ref{lem: hidden lemma}: if there are no 1-entries inside $M^X|_{I',J'}$, then during the \greedy execution all elements in the interval $[x+1, y-1]$ were $[t+1,t']$-hidden via $[x+1, y-1]$.
\end{proof}

To finish the proof, notice that if $T^X$ were to contain $(M_P \otimes G)$, then we could surround ``each copy'' of $G$ inside $(M_P \otimes G)$ with pairwise non-overlapping bounding boxes. Due to Lemma~\ref{forcing2}, each of these boxes must contain an access point, furthermore the relative positions of these access points is the same as that of the one-entries of $M_P$, contradicting that the input was $P$-free.
\fi
\subsection{$k$-increasing sequences}
\label{sub:k inc seq}
We prove Lemma~\ref{lem:input-revealing} (iii) and (iv).   

\begin{lemma}[Lemma~\ref{lem:input-revealing}(iii)]
$\inc_{k+1}$ ($\dec_{k+1})$ is a $k$-increasing ($k$-decreasing)
		gadget for \greedy where $\inc_{k}=(k,\dots,1)$ and \textup{$\dec_{k}=(1,2,\dots,k)$}. 
\end{lemma}

\begin{proof}

We only prove that $\inc_{k+1}$ is a $k$-increasing gadget. Let $q_{0},\dots,q_{k}$ be
	touch points in $\Greedy_{T}(X)$ such that $q_{i}.y<q_{i+1}.y$
	for all $i$ and form $\inc_{k+1}$ with bounding box $B$. 
	Note that, for all $1\le i\le k$, we can assume that there is no touch point $q'_i$ where $q'_i.x = q_i.x$ and $q_{i-1}.y<q'_i.y<q_i.y$, otherwise we can choose $q'_i$ instead of $q_i$.
	So, for any $t \in [q_{i-1}.y,q_i.y)$ we have that $\tau(q_{i-1}.x,t)\ge\tau(q_{i}.x,t)$ 
	for all $1\le i\le k$. Suppose that there is no access point in $B$.
	
	By \Cref{prop:hidden}~(ii), $q_{1}.x$ is hidden in $[1,q_{0}.x)$ after $q_0.y$. 
	So there must be an access point $p_1 \in[1,B.\xmin)\times(q_{0}.y,q_{1}.y]$, otherwise $q_1$ cannot be touched.
	We choose such $p_{1}$ where $p_{1}.x$ is maximized. If $p_{1}.y<q_{1}.y$,
	then $\tau(p_{1}.x,p_{1}.y),\tau(q_{0}.x,p_{1}.y)\ge\tau(q_{1}.x,p_{1}.y)$.
	So by \Cref{prop:hidden}~(iii), $q_{1}.x$ is hidden in $(p_{1}.x,q_{0}.x)$ after $p_1.y$ and hence $q_{1}$ cannot be touched, which is a contradiction.
	So we have $p_{1}.y=q_{1}.y$ and $p_{1}.x<B.\xmin$.
	
	Next, we prove by induction from $i=2$ to $k$ that, given the
	access point $p_{i-1}$ where $p_{i-1}.y=q_{i-1}.y$ and $p_{i-1}.x<B.\xmin$,
	there must be an access point $p_{i}$ where $p_{i}.y=q_{i}.y$ and
	$p_{i-1}.x<p_{i}.x<B.\xmin$. Again, by \Cref{prop:hidden}~(iii), there
	must be an access point $p_{i}\in(p_{i-1}.x,q_{i-1}.x)\times(q_{i-1}.y,q_{i}.y]$
	otherwise $q_{i}$ cannot be touched. We choose the point $p_{i}$
	where $p_{i}.x$ is maximized. If $p_{i}.y<q_{i}.y$, then $\tau(p_{i}.x,p_{i}.y),\tau(q_{i-1}.x,p_{i}.y)\ge\tau(q_{i}.x,p_{i}.y)$.
	So by \Cref{prop:hidden}~(iii), $q_{i}.x$ is hidden in $(p_{i}.x,q_{i-1}.x)$ after $p_i.y$ and hence $q_{i}$ cannot be touched.
	So we have $p_{i}.y=q_{i}.y$ and $p_{i-1}.x<p_{i}.x<B.\xmin$. Therefore,
	we get $p_{1}.x<\dots<p_{k}.x<B.\xmin$ which concludes the proof.
\end{proof}

\begin{lemma}[Lemma~\ref{lem:input-revealing}(iv)]
 $\dec_{k+1}$ ($\inc_{k+1})$ is a capture gadget for \greedy that
		runs on $k$-increasing $(k$-decreasing$)$ sequence.
\end{lemma}
\begin{proof}
We only prove that $\inc_{k+1}$ is a capture gadget for \greedy
	running on $k$-decreasing sequence $X$. Since $\inc_{k+1}$ is an
	increasing gadget, if $\inc_{k+1}$ appears in $\Greedy_{T}(X)$ with
	bounding box $B$, then either there is an access point in $B$ or
	there are access points forming $(1,2,\dots,k)$. The latter case
	cannot happen because $X$ avoids $(1,2,\dots,k)$.
	So $\inc_{k+1}$ is a capture gadget.
\end{proof}	
	\iffalse
\subsection{Signed Greedy Result}

We can restrict {\sc Greedy} in a way that it satisfies only rectangles of one of the two types (.. and ..). The two outputs thus obtained might not be a feasible {\sc MinASS} solution. However, the sum of the two output-sizes (plus the input size) is a constant-approximation for the \emph{independent rectangle lower bound}, a bound that subsumes all other known lower bounds for $\opt$~\cite{DHIKP09, Harmon}. 

\begin{theorem} 
The cost of {\sc SignedGreedy} on $X = \pset_S$ is $O(m \cdot g(k))$, for a fixed function $g$ (to be specified later).
\end{theorem} 
\begin{proof} 
We prove the theorem in a similar way. Let $\textsc{Greedy}_{+}$ and $\textsc{Greedy}_{-}$ denote the variants of {\sc SignedGreedy} that satisfy rectangles of type (..), resp.\ (..), and let $T^X_{+}$, resp.\ $T^X_{-}$ be the corresponding touch matrices. Then {\sc SignedGreedy}$(X) = w(T^X_{+}) + w(T^X_{-}) + |X|$.


We define the forcing gadgets $G_{+} = \left( \begin{matrix}
  0 & 1 \\
  1 & 0 
 \end{matrix}\right)$, and $G_{-} = \left( \begin{matrix}
  1 & 0 \\
  0 & 1 
 \end{matrix}\right)$. We make the following claim.
 \begin{lemma}
 \label{lem:forcing_signed}
 If $X$ avoids $P$, then $T^X_{+}$ is $(M_P \otimes G_{+})$-free, and $T^X_{-}$ is $(M_P \otimes G_{-})$-free.
 \end{lemma}
 
The proof is analogous (but simpler) to the proof of Lemma~\ref{lem:forcing1}. This implies Theorem~\ref{thm:pseudo}(ii) via Lemma~\ref{lem: forbidden}(ii), and by the observation that if $P$ is a permutation of size $k$, then $M_P \otimes G_i$ is a $2k$-by-$2k$ permutation matrix, for $i \in \{+,-\}$. 
\end{proof}  

\begin{proof}
Consider a pattern $\sigma$ and name four points that appear in $M$ by $(a,t_a), (b,t_b), (c,t_c)$ and $(d, t_d)$ respectively such that $t_a < t_d< t_b < t_c$ (so we have $a < c < d <b$.) 
Assume for contradiction that during the time interval $I = (s_I, t_I)$, entries in $J$ are never accessed.    

Define the time points $t'_a, t'_b, t'_c, t'_d$ before time interval $I$ for which $a,b,c,d$ are touched for the last time before entering $I$. 
In our earlier notation, $t'_a, t'_b, t'_c, t'_d$ are  $\tau(a, s_I), \tau(b,s_I), \tau(c,s_I), \tau(d,s_I)$, respectively.  


\begin{lemma} 
Let $q$ be any elements in $J$ such that there are two other elements $p$ and $r$ where $p < q <r$ and $t_q > \max (t_p,t_r)$. Then  
$t'_q > \max_{z \in J: z \neq q} t'_z$. 
\end{lemma} 
\begin{proof} 
There are two cases to consider. 
If both $t'_p$ and  $t'_r$ are at least $t'_q$, we are immediately done because $q$ is $(J, I)$-hidden and so would have never been touched at time $t'_q$. 


 
Now assume that $t'_p < t'_q \leq t'_r$ (the other case when $t'_r < t'_q \leq t'_p$ is very similar). 
Let $\tilde t$ be the first time we access on the left-side of $J$ (this time must exist before $t_p$ because otherwise $p$ is never touched. 
First, we claim that $\tau(q, \tilde t+1) \leq \tau(r',\tilde t+1)$ for some $r': q<r' \leq r$. This is because, from time $s_I$ to $\tilde t -1$, the element $q$ is never touched ($r$ blocks it).   
At time $\tilde t$, if $q$ is touched, some $q < r' \leq r$ would have been touched too ($r'$ could be the same as $r$).

Next, we claim also that $\tau(q, \tilde t+1) \leq \tau(p', \tilde t +1)$ for some $p': p \leq p' <q$. This is because $\tau(p, s_I) < \tau(q,s_I)$, and during $[s_I, \tilde t - 1]$, no access was from the left side. 
So we have maintained $\tau(p, \tilde t) < \tau(q, \tilde t)$.   
This means that either $p$ or some other $p'$ on the left of $p$ is in $stair_{\tilde t}(x_{\tilde t})$ and would have been touched.  
The previous reasoning implies that $q$ is hidden.  

\end{proof} 
\lk{we should define "hidden" in the preliminary}\ \\
We now apply the lemma twice: First on $\{a, c, b\}$ to say that $t'_c$ is the (strictly) maximum among $\{t'_z\}$, and then on $\{a,d,b\}$ to say that $t'_d$ is also the maximum, a contradiction.  
\end{proof} 


---------------------------------------------------------------   Kurt ---------------------------------------------------

I do not understand this proof. I tried to do my own and got stuck. Consider the following scenario.

\[ \begin{array}{cccc}
a & c & d & b \\ \hline
0 & 1 & 0 & 0 \\
1 & 0 & 1 & 0 \\
0 & 0 & 0 & 1 \\
\end{array}\]

The table illustrates the situation before time $t_a$. The 1 indicates the time of the last access and time flows upwards. 

At time $t_a$ we access an item to the left of $a$ and obtain: 
\[ \begin{array}{cccc}
a & c & d & b \\ \hline
{\bf 1}& 1 & 0 & 0 \\
0 & 0 & 1 & 0 \\
0 & 0 & 0 & 1 \\
\end{array}\]
The top row now corresponds to time $t_a$. 

Between time $t_a$ and $t_b$ we access an element to the right of $b$. This element sees $c$ and $d$, but does not see $b$. We obtain. 

\[ \begin{array}{cccc}
a & c & d & b \\ \hline
0 & 1 & 1& 0 \\
1 & 0 & 0 & 0 \\
0 & 0 & 0 & 1 \\
\end{array}\]

The access at time $t_b$ sees $b$ and $d$. We obtain:
\[ \begin{array}{cccc}
a & c & d & b \\ \hline
0 & 0 & 1 & {\bf 1} \\
0 & 1 & 0 & 0 \\
1 & 0 & 0 & 0 \\
\end{array}\]
The next access is to the left of $a$. It sees $d$ and $c$ and maybe also $a$. We obtain:

\[ \begin{array}{cccc}
a & c & d & b \\ \hline
1 & {\bf 1} & 1 & 0 \\
0 & 0 & 0 & 1 \\  \end{array} \]
The next access is to the right of $b$. It sees $d$ and maybe also $b$. We obtain.

\[ \begin{array}{cccc}
a & c & d & b \\ \hline
0 & 0 & {\bf 1} & 1\\
1 & 1 & 0 & 0\\  \end{array} \]
The bold ones indicate the forbidden pattern. 

Next comes the standard view. Time flows up. I indicate the accessed element with a $*$ and the touched elements with a 1. The forbidden pattern is indicated in boldface.

\[\begin{array}{l | ccccccc}
& & a & c & d & b & & \\ \hline
t_d & 0 & 0 & 0 & {\bf 1} & 1 & * & 0  \\
t_c & * & 1 & {\bf 1}  & 1 & 0 & 0 & 0  \\
t_b & 0 & 0 & 0 & 1 & {\bf 1}& * & 1  \\
& 0 & 0 & 1 & 1 & 0  & 1 & *    \\
t_a &* & {\bf 1}& 1 & 0 & 0  & 0 & 0    \\
& 0 & 0 & 1 & 0 & 0 & 0 & 0 \\
& 0 & 1 & 0 & 1 & 0 & 0 & 0 \\
& 0 & 0 & 0 & 0 & 1 & 1 & 0 \\
\end{array}\]
Isn't this a counterexample to our proof? \bigskip

It is not necessarily a counterexample to the theorem. The counterexample crucially rests on the fact that the access at the time between $t_a$ and $t_b$ sees $c$ and $d$, but does not see $b$. I constructed this situation by putting a 1 right of $b$ in row corresponding to $\tau(b,t_a)$. A remedy might be to ask for another one in column $c$. Then we would need an access between times $t_b$ and $t_c$ that sees $c$ and $d$, but does not see $a$. No, it is not a remedy. Add one more column on the left and put an access into this column at a time between $t_b$ and $t_c$.The access will touch $b$ and $c$, but not $a$. 



\end{proof}

\begin{theorem}[Theorem 4.3, Forcing Gadget] Assume $G^X$ contains pattern $\sigma$. Let $I$ and $J$ be any minimal time and element intervals such that $G^X |_{I,J}$ contains $\sigma$. Then $M^X |_{I,J}$ contains a 1, i.e., some element of $J$ is accessed in time interval $I$. 
\end{theorem}
\begin{proof}
Let $(a,t_a)$, $(b,t_b)$, $(c,t_c)$, and $(d,t_d)$ be the four points in $G^X |_{I,J}$ that are required by $\sigma$ such that $a < c < d < b$, $t_a < t_b < t_c < t_d$, and $I = [t_a,t_d]$ and $J = [a,b]$. 

There is no touch in $[a] \times (t_a,t_b)$ and no touch in $(a,c) \times [t_a,t_b)$ as otherwise, we could use this element instead of $a$, contradicting minimality of $I$ and $J$. Similarly, there is no touch in $[b] \times (t_a,t_b)$ and no touch in $(d,b) \times (t_a, t_c)$. MORE PROPERTIES.
$c$ is not touched in $(t_b,t_c)$ and $d$ is not touched in $(t_c,t_d)$. 

If two non-consecutive points in $\sset{a,b,c,d}$ are touched at time $t_a$, then either $c$ or $d$ is $[t_a+1,t_d]$-hidden via $J$ and we have a contradiction. 

%%\marginpar{TODO: check definition of hidden; check definition of $\tau$. should it latest touch before $t$ or before or at $t$?}

So at time $t_a$, $a$ is touched and $c$ may be touched. $d$ and $b$ are not touched. If there is an $x \in (c,b]$ with $\tau(x,t_a) \ge \tau(c,t_a)$ or an $x \in (d,b]$ with $\tau(x,t_a) \ge \tau(d,t_a)$ then $c$ or $d$ is $I$-hidden via $J$, a contradiction. Thus $\tau(c,t_a) > \tau(x,t_a)$ for all $x \in (c,b]$ and $\tau(d,t_a) > \tau(x,t_a)$ for all $x \in (d,b]$. In other words, $\tau(x,t_a)$ is a non-increasing function of $x$ in $[a,b]$ with proper decreases after $c$ and $d$. Accesses left of $a$ maintain this invariant. 

Consider accesses right of $b$ before and up to $t_b$. They do not lead to a touch in $J$ left of $c$, by minimality. 
If $c$ is not touched, the invariant is maintained. If anything right of $c$ and in $J$ is touched, $c$ is also touched. 
If the first touch of $c$ is at $t_b$, then the last touch time of $c$, $d$, and $b$ becomes $t_b$ and $d$ is hidden afterwards, a contradiction. 

At $t_b$, $b$ is touched. In $(t_a,t_b)$, all touched elements of $J$ must be contained in $[c,d]$ by minimality. 

We must have $t_a = \tau(a,t_b) < \tau(c,t_b) < \tau(d,t_b) \le \tau(b,t_b)$ at time $t_b$ as otherwise either $c$ or $d$ is $(t_b,t_d]$-hidden via $J$.   \textcolor{red}{here I use $\tau(x,t) = $ time of latest touch at or before $t$}

Consider the first access $(y,t_y)$ after $t_a$ which causes a touch of $c$. Then $y > b$. After this access $\tau(x,\ )$ is a non-increasing function on $[c,b]$. This invariant is maintained by accesses left of $a$. 

Consider now the first access after $t_y$ and before and including $t_b$ that touches $d$. 
\end{proof}
\fi


\iffalse
\subsubsection{Initial Trees}

An initial tree is a structure of BST before the first access in access
sequence comes. In the tree view of BST execution, we know the structure
of the tree exactly in every time step, and hence the structure of
the initial tree as well. This is not clear, however, in geometric
view what is our initial tree. In fact, in all above discussion, the
initial tree is always \emph{split tree }\cite{DHIKP09} which is invisible
to us in geometric view because of its very property. In this section,
we show how to model any initial tree into the geometric view of BST
execution.

First, we need one simple fact. Given a binary search tree $T$, we
denote the search path of element $a$ in $T$ as $path(a,T)$. 
\begin{proposition}
\label{same tree=same paths}Two BSTs $T_{1}$ and $T_{2}$,
containing $[n]$, have the same structure iff $path(a,T_{1})=path(a,T_{2})$
as a set, for all $a\in[n]$.\end{proposition}
\begin{proof}
If $T_{1}$ and $T_{2}$ share the same structure, then clearly $path(a,T_{1})=path(a,T_{2})$
for any element $a$. Conversely, we prove by induction on the depth
of the trees. Let $r$ be the root of $T_{1}$, and hence the root
of $T_{2}$ as $path(r,T_{1})=path(r,T_{2})$. By induction hypothesis,
left subtrees of the root $r$ in $T_{1}$ and $T_{2}$ have the same
structure. Similarly for right subtrees of the root in $T_{1}$ and
$T_{2}$. Altogether, $T_{1}$ and $T_{2}$ have the same structure.
\end{proof}
The stair of element is, in some sense, the geometric view of the
search path. Indeed, $path(a,T)=\{b\mid b$ is touched if $a$ is
touched in $T\}$. Similarly, we know that $stair(a,t)=\{b\mid b$
is touched if $a$ is touched at time $t\}$. Thus, to model the initial
tree $T$ into our geometric view, by Proposition \ref{same tree=same paths},
it is enough make it such that $stair(a,1)=path(a,T)$.

We can achieve this in a very straightforward way: below row 1, we
put the stacks of points in each columns with different height, see
Figure \ref{fig:geo init tree}. The exact height of stack will not
matter; what matters is comparative height between stacks. Given a
tree $T$ rooted at $r$, then the height of stack at column $r$
is the highest among other columns corresponding to elements in $T$.
The height of stack in other columns is recursively described by the
left and right subtrees of $r$.
\begin{lemma}
With above construction, $stair(a,1)=path(a,T)$ for all $a\in[n]$.\end{lemma}
\begin{proof}
Induction.
\end{proof}
\fi
